% TODO: remove [ ] brackets around units

\subsection{Example: file download}
\label{introduction:download_example}

	Suppose you're downloading a $720$[MB] file from the internet to your computer.
	At $t=0$ you click ``Save as'' to start the download.
	Consider the function $f(t)$
	that describes the amount of disk space occupied by the partially-downloaded file.
	%	At time $t$,
	%	your browser reports the download progress as a percentage
	%	that corresponds to the fraction $\frac{f(t)}{720 \text{[MB]}}$.

	\subsubsection{Download rate}
	\label{introduction:download_rate}

		The \emph{derivative function} $f^{\prime\!}(t)$,
		pronounced ``\:\!\!$f$ prime,''
		describes how the function $f(t)$ changes over time.
		In our example $f^{\prime\!}(t)$ is the download speed.
		If your download speed is $f^{\prime\!}(t)=2$[MB/s],
		then the file size $f(t)$ will increase by 2[MB] each second.
		If you maintain this download speed,
		the file size will grow at a constant rate:
		$f(0)=0$[MB], $f(1)=2$[MB], $f(2)=4$[MB], $\ldots$, $f(100)=200$[MB],
		and so on until $t=360$[s] when the download will be done.

		To estimate the time that remains before the download completes,
		assuming the current speed stays roughly constant,
		we can divide the amount of data that remains to be downloaded
		by the current download speed:
		\[
			\text{time remaining at } t = \tfrac{ 720 - f(t) }{ f^{\prime\!}(t) } \quad [\textrm{s}].
		\]
		The bigger the derivative $f^{\prime\!}(t)$,
		the faster the download will finish.
		If your internet connection were 10 times faster,
		the download would finish 10 times more quickly.


	\subsubsection{Inverse problem}
	\label{introduction:inverse_problem}

		Let's now consider the download scenario
		from the point of view of the modem that connects your computer to the internet.
		Any data you download comes through the modem,
		so the modem knows the download rate $f^{\prime\!}(t)$[MB/s]
		at all times during the download.
		%
		The modem is separate from your computer,
		so it doesn't know the file size $f(t)$.
		Nevertheless,
		the modem can infer the file size at time $t$
		from the transmission rate $f^{\prime\!}(t)$.
		Think about it---if the modem sees data flowing through at the rate of $f^{\prime\!}(t)=2$[MB/s],
		then it knows that the data accumulated on your computer
		is growing at the rate of $2$[MB] each second.
		In calculus,
		we describe the total file size accumulated until time $t=\tau$ (the Greek letter \emph{tau})
		as the \emph{integral} of the download rate $f^{\prime\!}(t)$ between $t=0$ and $t=\tau$:
		\[
			f(\tau) \; = \; \int_{t=0}^{t=\tau}  f^{\prime\!}(t)\, dt.
		\]
		The symbol $\int$ is an elongated $S$ that stands for \emph{sum}.
		Indeed,
		the ``integral of $f^{\prime\!}(t)$ between $0$ and $\tau$''
		is in some sense the sum of $f^{\prime\!}(t)$
		during each time instant $dt$ between $t=0$ and $t=\tau$.
		To calculate the total accumulated file size,
		we split the time interval between $t=0$ and $t=\tau$
		into many short time intervals $dt$ of length $1$[s].
		During each second,
		the file size grows by $f^{\prime\!}(t)\,dt$,
		where $f^{\prime\!}(t)$ is the download rate measured in [MB/s],
		and $dt$ is a time interval of duration $1$[s].
		%	Note the units work out:
		%	the data downloaded during one second is $f^{\prime\!}(t)dt$[MB].
		%	The file size on your computer at $t=\tau$
		%	is the sum of these 1-second contributions $f^{\prime\!}(t)\,dt$
		%	as $t$ varies from $t=0$ to $t=\tau$.
		% during each second from $t=0$ until $t=\tau$.
		% The file size at time $t=\tau$ is equal to the sum of the data downloaded 

		%	corresponds to the total amount of downloaded data stored on your computer.
		%	During this download period,
		%	the change in file size is described by the integral
